%! Setting Up the SVD — Unit 7, Lesson 7.0 — Lesson Plan
\documentclass[10pt]{article}
\usepackage{linalg-article}
\usepackage{linalg-boxes}
\usepackage{graphicx}   % -article does NOT load graphicx; needed for thumbnails/screenshots
\graphicspath{{images/}}
\newcommand{\TallMath}[1]{$\displaystyle #1\rule[-1.4em]{0pt}{3.2em}$}
\newcommand{\uu}{\mathbf{u}}
\newcommand{\vv}{\mathbf{v}}
\newcommand{\ww}{\mathbf{w}}
\newcommand{\ee}{\mathbf{e}}
\newcommand{\zero}{\mathbf{0}}
\newcommand{\T}{^{\top}}
\newcommand{\UnitNumberName}{Unit 7: The Singular Value Decomposition \quad}
\newcommand{\LessonNumberName}{Lesson 7.0: Setting Up the SVD}

\begin{document}
\begin{center}
    {\Large\bfseries \CourseName: \SchoolYear \\
    {\small\bfseries \UnitNumberName \LessonNumberName}}
\end{center}

\begin{tcolorbox}[colback=blush, colframe=burgundy, boxrule=0.9pt, arc=2mm,
    left=2mm, right=2mm, top=1.5mm, bottom=1.5mm]
    \textbf{Primary Objective:} Students can describe what a matrix does to the \emph{unit circle}:
    it stretches it into an \emph{ellipse} whose two half-widths are the \emph{singular values}
    $\sigma_1\ge\sigma_2\ge0$. They can find a perpendicular pair of input directions $\vv_1,\vv_2$
    that a matrix sends to a perpendicular pair of output directions $\uu_1,\uu_2$, satisfying
    $A\vv_i=\sigma_i\uu_i$, and read off the singular values as the output lengths. They can name the
    \emph{Singular Value Decomposition} $A=U\Sigma V\T$ (rotate--stretch--rotate) and explain \emph{why}
    it matters --- \emph{every} matrix (of any shape) has one and the singular values are always real and
    nonnegative, unlike eigenvectors (Unit~6). They can describe where Unit~7 is going: \emph{finding}
    the singular values/vectors from $A\T A$ (7.1), compressing images (7.2), principal component analysis
    (7.3), and the victory of orthogonality (7.4).
\end{tcolorbox}

\renewcommand{\arraystretch}{1.4}
\begin{skillbox}[Priority Ideas \& Skills]{goldbox}
\begin{tabularx}{\linewidth}{@{}X|X@{}}
\textbf{Priority Ideas \& Skills} & \textbf{Key Understandings (the ``why'')} \\[2pt]
\hline
\vspace{-6pt}
\begin{itemize}[leftmargin=1.1em, nosep, after=\vspace{2pt}]
  \item A matrix maps the \textbf{unit circle} to an \textbf{ellipse}; read the \textbf{singular values} as its half-widths.
  \item State $A\vv_i=\sigma_i\uu_i$: a perpendicular input pair $\to$ a perpendicular output pair.
  \item Read $\sigma_1\ge\sigma_2\ge0$ off diagonal, swap, and symmetric matrices.
  \item Name the SVD $A=U\Sigma V\T$ and why every matrix has one.
\end{itemize}
&
\vspace{-6pt}
\begin{itemize}[leftmargin=1.1em, nosep, after=\vspace{2pt}]
  \item Singular values measure \emph{how far} a matrix stretches the circle into an oval.
  \item ``Perpendicular in, perpendicular out'' works for \emph{any} matrix --- not just symmetric ones.
  \item Singular values are \emph{lengths}: always real, never negative (unlike eigenvalues).
  \item The SVD splits any matrix into rotate--stretch--rotate; it drives compression \& data science.
\end{itemize}
\end{tabularx}
\end{skillbox}

\begin{skillbox}[{Vocabulary, Concepts, \& Theorems}]{greenbox}
\renewcommand{\arraystretch}{1.3}
\begin{tabularx}{\linewidth}{@{}l X@{}}
\textbf{Singular value $\sigma$} & A half-width (semi-axis length) of the ellipse the unit circle maps to; the factor $\ge0$ a singular direction is stretched by. \\
\textbf{Input singular vectors $\vv_i$} & A perpendicular set of \emph{input} directions ($\vv_1\perp\vv_2$), the right singular vectors. \\
\textbf{Output singular vectors $\uu_i$} & The perpendicular \emph{output} directions ($\uu_1\perp\uu_2$) --- the ellipse's axes, the left singular vectors. \\
\textbf{The relation $A\vv_i=\sigma_i\uu_i$} & $A$ sends input direction $\vv_i$ to $\sigma_i$ times output direction $\uu_i$: perpendicular in, perpendicular out. \\
\textbf{Singular Value Decomposition} & $A=U\Sigma V\T$: columns of $V,U$ are the (orthonormal) input/output singular vectors; $\Sigma$ is diagonal with the $\sigma$'s. \\
\textbf{Finding $\sigma$} (preview 7.1) & $\sigma=\sqrt{\lambda}$ for the eigenvalues $\lambda$ of the symmetric matrix $A\T A$. \\
\end{tabularx}
\end{skillbox}

\begin{fixedskillbox}[Activate Prior Knowledge \& Spiral Review]{blush}
    The Warm-Up rehearses the three ideas this lesson builds on:
    \textbf{(1)} computing a \emph{matrix times a vector} $A\vv$ (Unit~1) --- a matrix stretching a
    direction; \textbf{(2)} the \emph{dot-product perpendicular test} (Unit~4) --- $\vv\cdot\ww=0$ means a
    right angle, the tool for ``perpendicular in, perpendicular out''; and \textbf{(3)} the \emph{length}
    of a vector and \emph{normalizing} to a unit vector (Unit~4) --- singular vectors are reported as unit
    vectors. Debrief item~2 aloud: a dot product of $0$ is exactly the right-angle condition the SVD needs.
\end{fixedskillbox}

\begin{skillbox}[Hook]{blush}
    Show a matrix acting on the plane (Lesson~5.3), then focus on the \textbf{unit circle} --- all vectors
    of length~$1$. Ask: \textbf{``What shape does the circle become after we multiply every point by a
    matrix?''} Answer: an \emph{ellipse}. Name the payoff question of Unit~7: \textbf{``How far does the
    circle stretch, and in which perpendicular directions?''} Those stretch amounts are the
    \emph{singular values}; the perpendicular input/output directions are the \emph{singular vectors}. They
    run all of Unit~7 --- finding them (7.1), image compression (7.2), PCA (7.3), and the victory of
    orthogonality (7.4). Tie back to Unit~6: eigenvectors kept a direction \emph{fixed}; singular vectors
    send one perpendicular frame to \emph{another}, and they exist for \emph{every} matrix.
\end{skillbox}

\begin{skillbox}[Lesson]{blush}
\begin{multicols}{2}
\textbf{1. Circle $\to$ ellipse.} Multiplying every unit vector by a matrix turns the unit circle into an
ellipse. For a diagonal $D=\begin{bsmallmatrix} 3 & 0\\ 0 & 2 \end{bsmallmatrix}$, the axes stretch to
lengths $3$ and $2$: the \textbf{singular values} $\sigma_1=3,\ \sigma_2=2$ (larger first).

\textbf{2. Perpendicular in $\to$ perpendicular out.} The key fact: every matrix has unit vectors
$\vv_1\perp\vv_2$ and $\uu_1\perp\uu_2$ with $A\vv_i=\sigma_i\uu_i$. On the Unit~6 matrix
$A=\begin{bsmallmatrix} 2 & 1\\ 1 & 2 \end{bsmallmatrix}$: $A\begin{bsmallmatrix} 1\\1 \end{bsmallmatrix}=
3\begin{bsmallmatrix} 1\\1 \end{bsmallmatrix}$ and $A\begin{bsmallmatrix} 1\\-1 \end{bsmallmatrix}=
1\begin{bsmallmatrix} 1\\-1 \end{bsmallmatrix}$, so $\sigma_1=3,\ \sigma_2=1$; the ellipse is tilted $45^\circ$.

\columnbreak

\textbf{3. Input and output frames can differ.} For the swap $S=\begin{bsmallmatrix} 0 & 2\\ 1 & 0
\end{bsmallmatrix}$, $S\ee_1=\begin{bsmallmatrix} 0\\1 \end{bsmallmatrix}$, $S\ee_2=\begin{bsmallmatrix} 2\\0
\end{bsmallmatrix}$ --- perpendicular out, $\sigma_1=2,\ \sigma_2=1$, and the largest input $\vv_1=\ee_2$ is
sent to a \emph{different} output $\uu_1=\begin{bsmallmatrix} 1\\0 \end{bsmallmatrix}$.

\textbf{4. The SVD + why it matters.} Collect the pieces into $A=U\Sigma V\T$ (rotate--stretch--rotate).
It exists for \emph{every} matrix and $\sigma\ge0$ always --- unlike eigenvectors. \textbf{Finding them}
(7.1): $\sigma=\sqrt{\lambda}$ from the eigenvalues of $A\T A$. \textbf{Road ahead:} \textbf{7.1} finding
$\sigma$; \textbf{7.2} image compression; \textbf{7.3} PCA; \textbf{7.4} the victory of orthogonality.
\end{multicols}
\end{skillbox}

\begin{skillbox}[Explicit Instruction: Reading Singular Values Off a Matrix]{blush}
\begin{tabularx}{\linewidth}{@{}p{0.46\linewidth} X@{}}
\textbf{Steps (the read-off).}
\begin{enumerate}[leftmargin=1.3em, nosep, after=\vspace{2pt}]
  \item Pick a perpendicular input pair (the axes $\ee_1,\ee_2$, or given eigenvectors).
  \item Compute each output $A\vv_i$.
  \item Check the outputs are perpendicular: $(A\vv_1)\cdot(A\vv_2)=0$.
  \item The output \emph{lengths} are the singular values --- list the larger first ($\sigma_1\ge\sigma_2$).
  \item The output \emph{directions} $\uu_i=A\vv_i/\sigma_i$ are the output singular vectors.
\end{enumerate}
&
\textbf{Worked example} ($A=\begin{bsmallmatrix} 2 & 1\\ 1 & 2 \end{bsmallmatrix}$).
Inputs $\vv_1=\begin{bsmallmatrix} 1\\1 \end{bsmallmatrix}$, $\vv_2=\begin{bsmallmatrix} 1\\-1 \end{bsmallmatrix}$
($\vv_1\cdot\vv_2=0$). Outputs $A\vv_1=\begin{bsmallmatrix} 3\\3 \end{bsmallmatrix}$,
$A\vv_2=\begin{bsmallmatrix} 1\\-1 \end{bsmallmatrix}$; $(A\vv_1)\cdot(A\vv_2)=0$. Lengths (over $|\vv_i|$):
$\sigma_1=3,\ \sigma_2=1$. As unit vectors, divide by $\sqrt2$. (\textbf{7.1 preview:} these same
$\sigma$'s are $\sqrt{9},\sqrt{1}$ from the eigenvalues of $A\T A=\begin{bsmallmatrix} 5 & 4\\ 4 & 5 \end{bsmallmatrix}$.)
\end{tabularx}
\end{skillbox}

\begin{skillbox}[Active Monitoring]{redbox}
Circulate during the notes practice and Tier work. Watch for and correct:
\begin{itemize}[leftmargin=1.2em, nosep]
  \item arithmetic slips in $A\vv$ (row $\cdot$ vector) --- every read-off depends on the product;
  \item forgetting the ordering convention $\sigma_1\ge\sigma_2$ (list the \emph{larger} stretch first);
  \item reporting a singular value as \emph{negative} --- a singular value is a length, always $\ge0$;
  \item conflating the input direction $\vv_i$ with the output direction $\uu_i$ (they differ in general).
\end{itemize}
Cold-call prompts: ``What shape does the unit circle become?'' ``Which output is longer --- so which is
$\sigma_1$?'' ``Are the two outputs perpendicular?'' ``Is $\uu_1$ the same as $\vv_1$ here?''
\end{skillbox}

\begin{skillbox}[Group Work \& Differentiation]{redbox}
\textbf{Stretching the Unit Circle into an Ellipse} activity (tiered; mirrors the notes).
\begin{multicols}{3}
\textbf{Tier R --- Remediate.} Diagonal matrices: compute $A\ee_1,A\ee_2$, describe the ellipse, read
$\sigma_1,\sigma_2$, and practice the ordering convention (larger first).

\columnbreak
\textbf{Tier A --- Approaching.} A swap matrix showing input frame $\ne$ output frame; then a symmetric
matrix (Unit~6) whose perpendicular eigenvectors \emph{are} the singular vectors and whose positive
eigenvalues \emph{are} the singular values; confirm ``perpendicular out.''

\columnbreak
\textbf{Tier E --- Extension.} \emph{First look at finding singular values:} compute $A\T A$, find its
eigenvalues, and take $\sigma=\sqrt{\lambda}$ --- recovering the notes' $\sigma=3,1$. This is the
Lesson~7.1 method.
\end{multicols}
\end{skillbox}

\begin{skillbox}[Individual Work \& Assessment]{redbox}
\textbf{Exit Ticket} (independent, no notes): read the singular values off a diagonal matrix (describe
the ellipse); compute $S\ee_1,S\ee_2$ for a swap matrix, check ``perpendicular out,'' and give
$\sigma_1,\sigma_2$; and a \textbf{conceptual/justification check} --- state in one sentence what a
singular value tells you geometrically. Collect and sort into ``got it / arithmetic slip / concept slip''
to decide whether 7.1 opens with a quick ``circle $\to$ ellipse'' re-cap.
\end{skillbox}

\begin{skillbox}[Reinforcement \& Extension]{goldbox}
\textbf{Homework:} read $\sigma$'s off diagonal matrices (incl.\ the ordering and the circle-stays-a-circle
case); a swap matrix separating input vs.\ output frames; a symmetric matrix tying singular values to
Unit~6 eigenvalues; and short justifications ($\sigma\ge0$ as a length; why every matrix has an SVD).
\textbf{Extension:} compute $A\T A$ and take $\sigma=\sqrt{\lambda}$ --- a first taste of Lesson~7.1.
\textbf{Preview:} Lesson~7.1, \emph{Singular Values and Singular Vectors} --- how to \emph{find} the
$\sigma$'s and singular vectors for \emph{any} matrix from $A\T A$, with $\uu_i=A\vv_i/\sigma_i$.
\end{skillbox}
\end{document}
